\documentclass[pdflatex,sn-basic,Numbered]{sn-jnl}

\usepackage{graphicx}
\usepackage{amsmath,amssymb,amsfonts}
\usepackage{amsthm}
\usepackage{mathtools}
\usepackage{enumitem}

\theoremstyle{thmstyleone}
\newtheorem{theorem}{Theorem}[section]
\newtheorem{proposition}[theorem]{Proposition}
\newtheorem{corollary}[theorem]{Corollary}
\newtheorem{lemma}[theorem]{Lemma}
\theoremstyle{thmstyletwo}
\newtheorem{example}[theorem]{Example}
\newtheorem{remark}[theorem]{Remark}
\theoremstyle{thmstylethree}
\newtheorem{definition}[theorem]{Definition}

\raggedbottom

\newcommand{\RR}{\mathbb{R}}
\newcommand{\NN}{\mathbb{N}}
\newcommand{\EE}{\mathbb{E}}
\newcommand{\Var}{\mathrm{Var}}
\newcommand{\PP}{\mathbb{P}}
\newcommand{\KL}{\mathrm{KL}}

\begin{document}

\title[Tilt Dynamics of Cylinder Information]{Tilt Dynamics of Cylinder Information and the Parry Measure on the Golden-Mean Shift}

\author[1]{\fnm{Haobo} \sur{Ma}}\email{auric@aelf.io}

\author*[2]{\fnm{Wenlin} \sur{Zhang}}\email{e1327962@u.nus.edu}

\affil[1]{\orgname{AELF PTE LTD.},
  \orgaddress{\street{\#14-02, Marina Bay Financial Centre Tower 1, 8 Marina Blvd},
  \city{Singapore}, \postcode{018981}, \country{Singapore}}}

\affil[2]{\orgname{National University of Singapore},
  \orgaddress{\city{Singapore}, \country{Singapore}}}

\abstract{We study exponential tilts of cylinder information $\tau_n(x)=-\log\mu([x_0^{n-1}])$ in the one-step Markov family on the golden-mean shift. The tilted problem closes inside this family: after Doob normalization the tilted law is again a one-step Markov law~$\mu_{q_p(t)}$, and the parameter map satisfies $L(q_p(t))=(1-t)L(p)$ where $L(p)=\log((1-p)/p^2)$. We compute the large-deviation rate function in relative-entropy coordinates, identify the Parry measure as the unique zero-jitter law, and give an exact normal form decomposing $\tau_n-nh(p)$ into an occupation-number bulk term and an endpoint correction. The tilt structure extends to the free-energy level via a composition law $F_{q_p(t)}(s)=F_p(s+t-st)-(1-s)F_p(t)$, which yields a variance transfer formula and characterizes the Parry measure through the constancy of its R\'enyi entropy spectrum. For general mixing shifts of finite type, we establish a Parry-relative fluctuation normal form, prove that the Parry measure is the unique zero-jitter law in the full one-step Markov simplex, and show that the quadratic coefficient $\sigma^2_\tau\sim 2(\log\lambda-h)$ near the Parry point is universal.}

\keywords{Golden-mean shift, Cylinder information, Parry measure, Exponential tilts, Asymptotic variance, Shifts of finite type}

\pacs[MSC Classification]{37B10, 37A50, 37D35, 60F10}

\maketitle

\section{Introduction}

Entropy, whether topological or measure-theoretic, gives a single number summarizing the complexity of a dynamical system. A finer question concerns the fluctuation structure of the \emph{cylinder information}
\[
\tau_n(x):=-\log\mu([x_0^{n-1}])
\]
around its linear growth rate. The Shannon--McMillan--Breiman theorem guarantees that $\tau_n/n$ converges to the entropy rate for ergodic measures, and central limit theorems, large-deviation principles, and related results describe the fluctuations for broad classes of mixing measures; see Haydn~\cite{Haydn2012} and Chazottes and Ugalde~\cite{ChazottesUgalde2005}. A natural next question is whether the \emph{exponential tilts} of~$\tau_n$---the measures obtained by reweighting orbits proportionally to $e^{t\tau_n}$---can be described explicitly in some concrete symbolic family.

In this paper we carry out such a description for the golden-mean shift
\[
X_\phi:=\bigl\{x\in\{0,1\}^{\NN}: x_kx_{k+1}=0 \text{ for all } k\ge 0\bigr\}
\]
and the one-step Markov family $\mu_p$ on~$X_\phi$ (parametrized by the single free parameter $p=P_{00}\in(0,1)$). The Parry measure---the measure of maximal entropy, corresponding to $p=\phi^{-1}$---plays a distinguished role throughout.

The basic observation is that exponential tilting preserves the family: after Doob normalization, the tilted law is again $\mu_{q_p(t)}$ for an explicit parameter $q_p(t)$. In the coordinate $L(p):=\log\frac{1-p}{p^2}$, the tilt map satisfies $L(q_p(t))=(1-t)L(p)$, and is thus conjugate to the linear map $u\mapsto(1-t)u$. The unit tilt sends every parameter to the Parry point~$\phi^{-1}$, and for $t\in(0,2)$ the Parry point is a global attractor.

This leads to several consequences. The centered cylinder information admits an exact normal form (Theorem~\ref{thm:normal-form}) that decomposes $\tau_n-nh(p)$ into a bulk term proportional to the centered occupation number---with coefficient $L(p)$, which vanishes at the Parry point---and an endpoint correction. Theorem~\ref{thm:tilted-free-energy} gives the Doob-transform identification and an explicit scaled cumulant generating function. Corollary~\ref{cor:tilt-dynamics} shows that the tilt maps form a multiplicative semigroup, and Corollary~\ref{cor:tilted-ldp} expresses the large-deviation rate function in relative-entropy coordinates, recovering the identity $\log\phi-h(p)=D(\mu_p\|\mu_\phi)$ as a special case. For the Parry measure itself, $\tau_n-n\log\phi$ depends only on the two endpoints of the word (Theorem~\ref{thm:exact-clock}), which gives a three-point defect law and an exact bounded variance formula.

The semigroup structure extends from the parameter level to the free-energy level: the composition law $F_{q_p(t)}(s)=F_p(s+t-st)-(1-s)F_p(t)$ (Theorem~\ref{thm:composition}) shows that tilting intertwines the entire fluctuation spectrum. An immediate corollary is a variance transfer formula $\sigma^2(q_p(t))=(1-t)^2 F_p''(t)$ (Corollary~\ref{cor:variance-transfer}), and a characterization of the Parry measure as the unique member of the family whose R\'enyi entropy rate is constant across all orders.

While the general theory of large deviations for Markov additive functionals is well developed (see Kifer~\cite{Kifer1990} and Dembo and Zeitouni~\cite{DemboZeitouni2010}), what distinguishes the present setting is that the \emph{tilted measures stay in the same parametric family} and the resulting dynamics on the parameter interval can be linearized globally. We show that the tilt closure itself extends to every mixing shift of finite type (Proposition~\ref{prop:general-tilt}), but that linearization requires a one-dimensional parameter space---a condition the golden-mean shift satisfies naturally because the adjacency constraint $A_{11}=0$ fixes the second row of the transition matrix.

Beyond tilt closure, the golden-mean results lift to the full Markov simplex of any mixing SFT. Theorem~\ref{thm:general-normal-form} provides a Parry-relative fluctuation normal form valid for all one-step Markov measures, Theorem~\ref{thm:general-rigidity} proves Parry rigidity via a finite-dimensional Perron--Frobenius argument, and Theorem~\ref{thm:universal} establishes the universality of the quadratic coefficient~$2$ in $\sigma^2_\tau\sim 2(\log\lambda-h)$.

The connection to the broader Bowen--Ruelle setting is made in Theorem~\ref{thm:gibbs}: for H\"older equilibrium states on a mixing shift of finite type, zero asymptotic variance of~$\tau_n$ (``zero jitter'') is equivalent to the potential being cohomologous to a constant. This is a standard consequence of exponential decay of correlations for Gibbs measures \cite{Bowen1975,Ruelle1978,ParryPollicott1990,Baladi2000}, and the golden-mean calculations of Sections~3--4 provide a fully explicit instance.

The classical background comes from Parry's intrinsic Markov chains~\cite{Parry1964} and from the Bowen--Ruelle theory \cite{Bowen1975,Ruelle1978,ParryPollicott1990,Walters1982}. The martingale approach to variance asymptotics for stationary processes goes back to Gordin~\cite{Gordin1969}; see Baladi~\cite{Baladi2000} for a modern treatment via transfer operators. The information-theoretic perspective on cylinder measures appears in Shields~\cite{Shields1996} and Kontoyiannis and Verd\'u~\cite{KontoyiannisVerdu2014}. Large deviations for symbolic systems are further studied in Pollicott and Sharp~\cite{PollicottSharp2009}; intrinsic ergodicity beyond the specification property is treated by Climenhaga and Thompson~\cite{ClimenhagaThompson2012}. More recent discussions of distinguished symbolic measures include the natural-measure perspective of Hu and Lin~\cite{HuLin2013} and the characteristic-measure viewpoint of Frisch and Tamuz~\cite{FrischTamuz2022}.

The paper is organized as follows. We fix notation and recall the Parry measure in Section~2. Section~3 contains the exact finite-time formulas at the Parry point. Section~4 develops the full-parameter normal form, the tilt dynamics, the free-energy composition law, large deviations, and the entropy-deficit identities---this is the heart of the paper. In Section~5 we extend the theory to general mixing SFTs: tilt closure, the Parry-relative normal form, Parry rigidity in the Markov simplex, and tangential linearization. Section~6 treats H\"older equilibrium states, the second-order comparison near the Parry point, and the universality of the quadratic coefficient~$2$. We close with open questions in Section~7.

\section{Preliminaries}

Let
\[
A=\begin{pmatrix}1&1\\1&0\end{pmatrix}
\]
be the adjacency matrix of the golden-mean shift. Its Perron eigenvalue is the golden ratio
\[
\phi:=\frac{1+\sqrt{5}}{2},
\]
with right and left Perron eigenvectors
\[
r=(\phi,1)^\top,
\qquad
\ell=(1,\phi^{-1})^\top,
\qquad
\ell^\top r=\sqrt{5}.
\]

\begin{definition}[Cylinder information and asymptotic jitter]
Let $\mu$ be a stationary law on $X_\phi$. For $x\in X_\phi$ and $n\ge 1$, define
\[
\tau_n(x):=-\log\mu([x_0^{n-1}]).
\]
If the limit exists, the asymptotic jitter rate is
\[
\sigma_\tau^2(\mu):=\lim_{n\to\infty}\frac{1}{n}\Var_\mu(\tau_n).
\]
We say that $\mu$ has \emph{zero jitter} when $\sigma_\tau^2(\mu)=0$. (The term ``jitter'' is used as a shorthand for the asymptotic variance rate of cylinder information; it is not standard terminology.)
\end{definition}

The Parry measure $\mu_\phi$ on $X_\phi$ is the stationary Markov law with transition matrix
\[
P_\phi=
\begin{pmatrix}
\phi^{-1} & \phi^{-2} \\
1 & 0
\end{pmatrix}
\]
and stationary distribution
\[
\pi_\phi=
\left(\frac{\phi}{\sqrt{5}},\frac{1}{\phi\sqrt{5}}\right)
=
\left(\frac{5+\sqrt{5}}{10},\frac{5-\sqrt{5}}{10}\right).
\]

More generally, every ergodic one-step Markov law on $X_\phi$ is parametrized by
\[
P_p=
\begin{pmatrix}
p & 1-p \\
1 & 0
\end{pmatrix},
\qquad 0<p<1,
\]
with stationary vector
\[
\pi_p=\left(\frac{1}{2-p},\frac{1-p}{2-p}\right).
\]
Its entropy rate is
\[
h(p)=\frac{H_2(p)}{2-p},
\qquad
H_2(p):=-p\log p-(1-p)\log(1-p).
\]

\section{The Parry Measure}

We begin with exact formulas for cylinder information under the Parry measure. The key point is that $\tau_n$ decomposes into a linear term and a correction that depends only on the first and last symbol.

\begin{theorem}[Parry cylinder formula]\label{thm:exact-clock}
Let $\mu_\phi$ be the Parry measure on $X_\phi$. For every legal word $a_{0:n-1}\in X_\phi$,
\begin{equation}\label{eq:exact-clock}
\tau_n(a_{0:n-1})
=
(n-1)\log\phi+\log\sqrt{5}-\log(\ell_{a_0}r_{a_{n-1}}).
\end{equation}
Equivalently,
\begin{equation}\label{eq:endpoint-trichotomy}
\tau_n(a_{0:n-1})-n\log\phi
=
\log\sqrt{5}-2\log\phi+(a_0+a_{n-1})\log\phi,
\end{equation}
so the defect from the linear term $n\log\phi$ depends only on the two endpoints of the word and takes exactly three values.
\end{theorem}

\begin{proof}
For any irreducible shift of finite type with $0$-$1$ adjacency matrix $A$, Perron eigenvalue $\lambda$, right eigenvector $r$, and left eigenvector $\ell$, the Parry measure of a legal cylinder satisfies
\[
\mu_P([a_0\cdots a_{n-1}])
=
\frac{\ell_{a_0}r_{a_{n-1}}}{\ell^\top r}\lambda^{-(n-1)}.
\]
Indeed, each legal transition contributes the factor $r_{a_{k+1}}/(\lambda r_{a_k})$, and the telescoping product leaves $r_{a_{n-1}}/(\lambda^{n-1}r_{a_0})$, which is multiplied by the initial weight $\ell_{a_0}r_{a_0}/(\ell^\top r)$. Specializing to the golden-mean shift gives \eqref{eq:exact-clock}. Since
\[
\ell_0r_0=\phi,\qquad \ell_0r_1=1,\qquad \ell_1r_0=1,\qquad \ell_1r_1=\phi^{-1},
\]
formula \eqref{eq:endpoint-trichotomy} follows immediately.
\end{proof}

\begin{corollary}[Endpoint law and finite-time variance]\label{cor:finite-time-law}
Let
\[
D_n:=\tau_n-n\log\phi.
\]
Under the Parry measure $\mu_\phi$, $D_n$ is supported on the three values
\[
d_-:=\log\sqrt{5}-2\log\phi,\qquad
d_0:=\log\sqrt{5}-\log\phi,\qquad
d_+:=\log\sqrt{5},
\]
with exact probabilities
\begin{align}
\PP_{\mu_\phi}(D_n=d_-)&=\left(\frac{5+\sqrt{5}}{10}\right)^2+\frac{1}{5}(-\phi^{-2})^{n-1}, \label{eq:defect-law-minus}\\
\PP_{\mu_\phi}(D_n=d_0)&=\frac{2}{5}\bigl(1-(-\phi^{-2})^{n-1}\bigr), \label{eq:defect-law-zero}\\
\PP_{\mu_\phi}(D_n=d_+)&=\left(\frac{5-\sqrt{5}}{10}\right)^2+\frac{1}{5}(-\phi^{-2})^{n-1}. \label{eq:defect-law-plus}
\end{align}
Consequently,
\begin{equation}\label{eq:finite-time-variance}
\Var_{\mu_\phi}(\tau_n)
=
\frac{2}{5}\bigl(1+(-\phi^{-2})^{n-1}\bigr)(\log\phi)^2.
\end{equation}
In particular, the Parry measure has bounded finite-time variance for cylinder information and
\[
\Var_{\mu_\phi}(\tau_n)\longrightarrow \frac{2}{5}(\log\phi)^2.
\]
\end{corollary}

\begin{proof}
By \eqref{eq:endpoint-trichotomy},
\[
D_n=d_-+(X_0+X_{n-1})\log\phi.
\]
Hence the three defect values correspond to the endpoint sums $X_0+X_{n-1}\in\{0,1,2\}$. Since $\mu_\phi$ is stationary Markov,
\[
\PP(X_0=i,X_{n-1}=j)=\pi_{\phi,i}(P_\phi^{n-1})_{ij}.
\]
The second eigenvalue of $P_\phi$ is $-\phi^{-2}$, and a direct two-state computation gives
\[
(P_\phi^m)_{11}=\pi_{\phi,1}+\pi_{\phi,0}(-\phi^{-2})^m,
\qquad
(P_\phi^m)_{01}=\pi_{\phi,1}-\pi_{\phi,1}(-\phi^{-2})^m,
\]
with analogous formulas for the remaining entries. Substituting $m=n-1$, together with
\[
\pi_{\phi,0}\pi_{\phi,1}=\frac{1}{5},
\]
yields \eqref{eq:defect-law-minus}--\eqref{eq:defect-law-plus}.

For the variance, write
\[
D_n=d_-+(I_0+I_{n-1})\log\phi,
\qquad
I_k:=\mathbf{1}_{\{X_k=1\}}.
\]
Therefore
\[
\Var_{\mu_\phi}(\tau_n)=\Var_{\mu_\phi}(D_n)
=(\log\phi)^2\Var(I_0+I_{n-1}).
\]
Because $\Var(I_0)=\pi_{\phi,0}\pi_{\phi,1}=1/5$ and
\[
\mathrm{Cov}(I_0,I_{n-1})
=
\pi_{\phi,1}(P_\phi^{n-1})_{11}-\pi_{\phi,1}^2
=
\frac{1}{5}(-\phi^{-2})^{n-1},
\]
we obtain
\[
\Var(I_0+I_{n-1})
=
2\Var(I_0)+2\mathrm{Cov}(I_0,I_{n-1})
=
\frac{2}{5}\bigl(1+(-\phi^{-2})^{n-1}\bigr),
\]
which is exactly \eqref{eq:finite-time-variance}.
\end{proof}

\begin{remark}
The Parry measure is therefore more rigid than the asymptotic definition of zero jitter alone suggests. Not only is $\Var(\tau_n)=o(n)$, but the entire finite-time defect law is explicit and its transient oscillation is governed by the subleading eigenvalue $-\phi^{-2}$ of the Parry transition matrix.
\end{remark}

\section{The One-Step Markov Family}

We now turn to the full one-step Markov family $\{\mu_p : 0<p<1\}$ on~$X_\phi$. Our starting point is an exact decomposition of the centered cylinder information that holds for every parameter~$p$.

\begin{theorem}[Full-parameter information normal form]\label{thm:normal-form}
For every $p\in(0,1)$, $n\ge 1$, and admissible word $x_0^{n-1}\in X_\phi$, let $T_n(x):=\sum_{k=0}^{n-1}x_k$ denote the occupation number of symbol~$1$. Then
\begin{equation}\label{eq:normal-form}
\tau_n(x) - nh(p) = -L(p)\!\left(T_n(x)-\frac{n(1-p)}{2-p}\right)+B_p(x_0,x_{n-1}),
\end{equation}
where $B_p(i,j):=\log(2-p)+(1-i-j)\log p$. In particular, the linear term in $T_n$ vanishes if and only if $L(p)=0$, i.e., $p=\phi^{-1}$, recovering the endpoint defect law of Theorem~\ref{thm:exact-clock} as the degeneration of \eqref{eq:normal-form} at the Parry point.
\end{theorem}

\begin{proof}
In the golden-mean shift, each symbol~$1$ forces a subsequent~$0$, so the edge counts in an admissible word of length~$n$ satisfy
\[
N_{10}=T_n-x_{n-1},\quad N_{01}=T_n-x_0,\quad N_{00}=n-1-2T_n+x_0+x_{n-1}.
\]
Since $\tau_n=-\log\pi_p(x_0)-N_{00}\log p-N_{01}\log(1-p)$ and $-\log\pi_p(x_0)=\log(2-p)-x_0\log(1-p)$, the $x_0\log(1-p)$ terms cancel and
\[
\tau_n=\log(2-p)-(n-1)\log p-L(p)\,T_n-(x_0+x_{n-1})\log p,
\]
where we used $2\log p-\log(1-p)=-L(p)$. On the other hand, $h(p)=-\log p-(1-p)L(p)/(2-p)$ (a direct rewriting of $h=H_2/(2-p)$), so $nh(p)=-n\log p-n(1-p)L(p)/(2-p)$. Subtracting gives \eqref{eq:normal-form}.
\end{proof}

\begin{remark}\label{rem:normal-form}
The decomposition \eqref{eq:normal-form} separates the fluctuations of $\tau_n$ into a bulk term proportional to the centered occupation number $T_n-\EE_{\mu_p}T_n$ and a boundary correction that depends only on the two endpoints. The bulk coefficient is~$L(p)$, which vanishes precisely at the Parry point. This explains \emph{structurally} why the Parry measure has bounded variance: its cylinder information is insensitive to the interior of the word.
\end{remark}

We now derive a general criterion for the vanishing of asymptotic variance, based on the Poisson equation.

\begin{lemma}[Zero-variance criterion for finite-state Markov additive functionals]\label{lem:markov-zero-variance}
Let $(X_n)_{n\ge 0}$ be an ergodic finite-state Markov chain with transition matrix $P$, stationary distribution $\pi$, and edge observable $g(i,j)$ on allowed transitions. Define
\[
S_n:=\sum_{k=0}^{n-2}g(X_k,X_{k+1}),
\qquad
h:=\EE_\pi[g(X_0,X_1)],
\qquad
\bar g(i):=\sum_jP(i,j)g(i,j).
\]
Let $v$ solve the Poisson equation
\[
(P-I)v=h\mathbf{1}-\bar g.
\]
Then
\[
S_n-(n-1)h=M_n+v(X_0)-v(X_{n-1}),
\]
where $(M_n)$ is a martingale with stationary increments
\[
m(i,j):=g(i,j)-h+v(j)-v(i).
\]
Consequently,
\[
\lim_{n\to\infty}\frac{1}{n}\Var(S_n)=\EE_\pi[m(X_0,X_1)^2],
\]
and this limit vanishes if and only if
\[
g(i,j)=h+u(i)-u(j)
\]
for every allowed transition $(i,j)$, for some function $u$ on the state space.
\end{lemma}

\begin{proof}
The Poisson equation gives
\[
g(i,j)-h=m(i,j)+v(i)-v(j),
\]
so summing along a path yields the stated decomposition. Since the increments $m(X_k,X_{k+1})$ form a stationary square-integrable martingale difference sequence,
\[
\Var(M_n)=\sum_{k=0}^{n-2}\EE[m(X_k,X_{k+1})^2]=(n-1)\EE_\pi[m(X_0,X_1)^2].
\]
The boundary term $v(X_0)-v(X_{n-1})$ is uniformly bounded, so dividing by $n$ gives the asymptotic variance formula. The limit is zero if and only if $m(X_0,X_1)=0$ almost surely, which is equivalent to
\[
g(i,j)=h+v(i)-v(j)
\]
on every allowed edge.
\end{proof}

\begin{theorem}[Zero jitter in the one-step Markov family]\label{thm:parry-rigidity}
Among ergodic one-step Markov laws on the golden-mean shift, the Parry measure is the unique zero-jitter measure:
\[
\sigma_\tau^2(\mu_p)=0
\quad\Longleftrightarrow\quad
p=\phi^{-1}.
\]
For every $p\neq \phi^{-1}$, the centered cylinder information satisfies the central limit theorem
\[
\frac{\tau_n-nh(p)}{\sqrt{n}}
\Rightarrow
\mathcal{N}(0,\sigma^2(p))
\]
with strictly positive variance.
\end{theorem}

\begin{proof}
Let $g_p(i,j):=-\log P_p(i,j)$ be the edge information function. The cylinder information under $\mu_p$ is
\[
\tau_n=-\log \pi_p(X_0)+\sum_{k=0}^{n-2}g_p(X_k,X_{k+1}).
\]
If $g_p$ is cohomologous to a constant, say
\[
g_p(i,j)=c+u(i)-u(j),
\]
then the Birkhoff sum telescopes and $\Var(\tau_n)=O(1)$, hence $\sigma_\tau^2(\mu_p)=0$.

Conversely, if $\sigma_\tau^2(\mu_p)=0$, then $\tau_n$ and
\[
\sum_{k=0}^{n-2}g_p(X_k,X_{k+1})
\]
differ only by the bounded term $-\log\pi_p(X_0)$, so they have the same asymptotic variance rate. Lemma~\ref{lem:markov-zero-variance} therefore implies that $g_p$ is cohomologous to a constant. Writing the coboundary equations for the three allowed transitions,
\[
-\log p=c,\qquad -\log(1-p)=c+u(0)-u(1),\qquad 0=c+u(1)-u(0),
\]
one obtains
\[
-\log(1-p)=-2\log p,
\]
equivalently
\[
1-p=p^2.
\]
The unique solution in $(0,1)$ is $p=\phi^{-1}$, namely the Parry measure.

For $p\neq \phi^{-1}$, the observable $g_p$ is not cohomologous to a constant. The standard central limit theorem for additive functionals of ergodic Markov chains \cite{Gordin1969} then yields the stated Gaussian fluctuations.
\end{proof}

\begin{proposition}[Asymptotic jitter formula]\label{prop:variance-formula}
For the Markov law $\mu_p$ on $X_\phi$,
\begin{equation}\label{eq:variance-formula}
\sigma^2(p)
=
\frac{p(1-p)}{(2-p)^3}
\left(\log\frac{1-p}{p^2}\right)^2.
\end{equation}
In particular, $\sigma^2(p)=0$ if and only if $p=\phi^{-1}$.
\end{proposition}

\begin{proof}
Let $S_n:=\sum_{k=0}^{n-2}g_p(X_k,X_{k+1})$. Since
\[
\tau_n=-\log\pi_p(X_0)+S_n,
\]
the difference between $\tau_n$ and $S_n$ is uniformly bounded in $n$, so they have the same asymptotic variance rate. By Lemma~\ref{lem:markov-zero-variance}, it therefore suffices to solve the Poisson equation
\[
(P_p-I)v=h(p)\mathbf{1}-\bar{g},
\qquad
\bar{g}(i):=\sum_jP_p(i,j)g_p(i,j).
\]
Fixing $v(1)=0$, one obtains $v(0)=h(p)$ and martingale increments
\[
m(0,0)=-\log p-h(p),\qquad
m(0,1)=-\log(1-p)-2h(p),\qquad
m(1,0)=0.
\]
Lemma~\ref{lem:markov-zero-variance} now yields
\[
\sigma^2(p)=\EE_{\pi_p}[m(X_0,X_1)^2]
=
\frac{1}{2-p}\left(p\,m(0,0)^2+(1-p)\,m(0,1)^2\right).
\]
Writing
\[
a:=-\log p,\qquad b:=-\log(1-p),
\]
and using $h(p)=(pa+(1-p)b)/(2-p)$, one finds
\[
a-h(p)=\frac{1-p}{2-p}(2a-b),
\qquad
b-2h(p)=-\frac{p}{2-p}(2a-b).
\]
Substituting these identities into the previous expression yields
\[
\sigma^2(p)
=
\frac{p(1-p)}{(2-p)^3}(2a-b)^2
=
\frac{p(1-p)}{(2-p)^3}\left(\log\frac{1-p}{p^2}\right)^2,
\]
which is \eqref{eq:variance-formula}.
\end{proof}

The next result is the heart of the paper. It shows that exponential tilting of cylinder information does not leave the one-step Markov family, and that the induced dynamics on the parameter interval are conjugate to linear multiplication.

\begin{theorem}[Tilted free energy and the Parry measure]\label{thm:tilted-free-energy}
For $p\in(0,1)$ and $t\in\RR$, define
\[
F_p(t):=\lim_{n\to\infty}\frac{1}{n}\log \EE_{\mu_p}\bigl[e^{t\tau_n}\bigr].
\]
Let
\[
B_p(t):=
\begin{pmatrix}
p^{1-t} & (1-p)^{1-t}\\
1 & 0
\end{pmatrix}.
\]
Then the following hold.
\begin{enumerate}[label=(\roman*),leftmargin=*,itemsep=2pt]
  \item The limit exists and
  \begin{equation}\label{eq:free-energy}
  F_p(t)=\log\rho_p(t),
  \qquad
  \rho_p(t):=
  \frac{p^{1-t}+\sqrt{p^{2(1-t)}+4(1-p)^{1-t}}}{2}.
  \end{equation}
  \item If
  \begin{equation}\label{eq:qpt-definition}
  q_p(t):=\frac{p^{1-t}}{\rho_p(t)},
  \end{equation}
  then the Doob transform of $B_p(t)$ is exactly $P_{q_p(t)}$. Equivalently,
  \begin{equation}\label{eq:defect-coordinate-flow}
  \frac{1-q_p(t)}{q_p(t)^2}
  =
  \left(\frac{1-p}{p^2}\right)^{1-t},
  \qquad\text{that is}\qquad
  L(q_p(t))=(1-t)L(p),
  \end{equation}
  where $L(p):=\log\frac{1-p}{p^2}$.
  \item The tilt maps $T_t(p):=q_p(t)$ satisfy
  \begin{equation}\label{eq:tilt-semigroup}
  T_s(T_t(p))=T_{s+t-st}(p)
  \end{equation}
  for all $s,t\in\RR$. In particular,
  \begin{equation}\label{eq:unit-tilt-parry}
  q_p(1)=\phi^{-1}
  \qquad\text{for every }p\in(0,1).
  \end{equation}
  \item The function $F_p$ is real analytic on $\RR$ and
  \begin{equation}\label{eq:Fprime-cross-entropy}
  F_p'(t)
  =
  \frac{-q_p(t)\log p-(1-q_p(t))\log(1-p)}{2-q_p(t)}
  =
  h(q_p(t))+D(\mu_{q_p(t)}\|\mu_p).
  \end{equation}
  Consequently,
  \begin{equation}\label{eq:relative-entropy-along-tilt}
  tF_p'(t)-F_p(t)=D(\mu_{q_p(t)}\|\mu_p).
  \end{equation}
\end{enumerate}
\end{theorem}

\begin{proof}
For a legal word $a_{0:n-1}$,
\[
\mu_p([a_0^{n-1}])=\pi_p(a_0)\prod_{k=0}^{n-2}P_p(a_k,a_{k+1}).
\]
Hence
\[
\EE_{\mu_p}[e^{t\tau_n}]
=
\sum_{a_0^{n-1}\in X_\phi}\mu_p([a_0^{n-1}])^{1-t}
=
\alpha_p(t)\,B_p(t)^{n-1}\mathbf{1},
\]
where
\[
\alpha_p(t):=(\pi_{p,0}^{1-t},\pi_{p,1}^{1-t}),
\]

The matrix $B_p(t)$ is primitive, so Perron--Frobenius yields
\[
\lim_{n\to\infty}\frac{1}{n}\log \EE_{\mu_p}[e^{t\tau_n}]
=
\log\rho(B_p(t)).
\]
The Perron eigenvalue of $B_p(t)$ is the positive root of
\[
\lambda^2-p^{1-t}\lambda-(1-p)^{1-t}=0,
\]
which is exactly $\rho_p(t)$ in \eqref{eq:free-energy}. This proves (i).

For (ii), note that $r_t:=(\rho_p(t),1)^\top$ is a positive right Perron eigenvector of $B_p(t)$. The normalized tilt is therefore the Markov kernel
\[
Q_t(i,j):=\frac{B_p(t)_{ij}\,r_t(j)}{\rho_p(t)\,r_t(i)}.
\]
A direct computation gives
\[
Q_t=
\begin{pmatrix}
\dfrac{p^{1-t}}{\rho_p(t)} & \dfrac{(1-p)^{1-t}}{\rho_p(t)^2}\\[1.5ex]
1 & 0
\end{pmatrix}
=
\begin{pmatrix}
q_p(t) & 1-q_p(t)\\
1 & 0
\end{pmatrix}
=P_{q_p(t)},
\]
where the identity $1-q_p(t)=(1-p)^{1-t}/\rho_p(t)^2$ follows from the eigenvalue equation for $\rho_p(t)$. Consequently,
\[
\frac{1-q_p(t)}{q_p(t)^2}
=
\frac{(1-p)^{1-t}\rho_p(t)^{-2}}{p^{2(1-t)}\rho_p(t)^{-2}}
=
\left(\frac{1-p}{p^2}\right)^{1-t},
\]
which is \eqref{eq:defect-coordinate-flow}.

Since $L'(p)=-\frac{1}{1-p}-\frac{2}{p}<0$, the map $L:(0,1)\to\RR$ is strictly decreasing and therefore injective. Applying \eqref{eq:defect-coordinate-flow} twice gives
\[
L(T_s(T_t(p)))=(1-s)L(T_t(p))=(1-s)(1-t)L(p),
\]
while
\[
L(T_{s+t-st}(p))=(1-(s+t-st))L(p)=(1-s)(1-t)L(p).
\]
Injectivity of $L$ proves \eqref{eq:tilt-semigroup}. Setting $t=1$ in \eqref{eq:defect-coordinate-flow} yields
\[
\frac{1-q_p(1)}{q_p(1)^2}=1,
\]
so $q_p(1)$ is the unique solution of $1-q=q^2$ in $(0,1)$, namely $\phi^{-1}$. This proves (iii).

For (iv), differentiate the Perron equation
\[
\rho_p(t)^2-p^{1-t}\rho_p(t)-(1-p)^{1-t}=0.
\]
Since
\[
\frac{d}{dt}p^{1-t}=-p^{1-t}\log p,
\qquad
\frac{d}{dt}(1-p)^{1-t}=-(1-p)^{1-t}\log(1-p),
\]
we obtain
\[
\rho_p'(t)\bigl(2\rho_p(t)-p^{1-t}\bigr)
=
-p^{1-t}\rho_p(t)\log p-(1-p)^{1-t}\log(1-p).
\]
Dividing by $\rho_p(t)^2$ and using
\[
\frac{p^{1-t}}{\rho_p(t)}=q_p(t),
\qquad
\frac{(1-p)^{1-t}}{\rho_p(t)^2}=1-q_p(t),
\qquad
\frac{2\rho_p(t)-p^{1-t}}{\rho_p(t)}=2-q_p(t),
\]
we find
\[
F_p'(t)=\frac{\rho_p'(t)}{\rho_p(t)}
=
\frac{-q_p(t)\log p-(1-q_p(t))\log(1-p)}{2-q_p(t)}.
\]
This is the cross-entropy rate of the Markov law $\mu_{q_p(t)}$ relative to the kernel $P_p$, so
\[
F_p'(t)=h(q_p(t))+D(\mu_{q_p(t)}\|\mu_p).
\]
Finally, for stationary one-step Markov laws on the same support,
\[
D(\mu_q\|\mu_p)
=
\frac{1}{2-q}
\left[
q\log\frac{q}{p}+(1-q)\log\frac{1-q}{1-p}
\right].
\]
With $q=q_p(t)$, the identities
\[
\log\frac{q_p(t)}{p}=-t\log p-F_p(t),
\qquad
\log\frac{1-q_p(t)}{1-p}=-t\log(1-p)-2F_p(t)
\]
follow from \eqref{eq:qpt-definition} and the formula for $1-q_p(t)$ above. Substituting them into the relative-entropy formula yields
\[
D(\mu_{q_p(t)}\|\mu_p)=tF_p'(t)-F_p(t),
\]
which is \eqref{eq:relative-entropy-along-tilt}.
\end{proof}

\begin{corollary}[Dynamics of the tilt maps]\label{cor:tilt-dynamics}
For each fixed $t\in\RR$, let $T_t:(0,1)\to(0,1)$ be given by $T_t(p)=q_p(t)$. Then:
\begin{enumerate}[label=(\roman*),leftmargin=*,itemsep=2pt]
  \item $L\circ T_t=(1-t)L$, so $T_t$ is conjugate, via the strictly decreasing coordinate $L$, to the linear map $u\mapsto (1-t)u$ on $\RR$.
  \item For every $n\ge 1$,
  \begin{equation}\label{eq:iterate-formula}
  T_t^{\,n}=T_{1-(1-t)^n}.
  \end{equation}
  \item If $0<t<2$, then $\phi^{-1}$ is the unique fixed point of $T_t$ and is globally attracting:
  \begin{equation}\label{eq:global-attractor}
  \bigl|L(T_t^{\,n}(p))\bigr|=|1-t|^n\,|L(p)|
  \longrightarrow 0
  \qquad\text{for every }p\in(0,1).
  \end{equation}
  In particular, $T_t^{\,n}(p)\to \phi^{-1}$ for every $p\in(0,1)$.
\end{enumerate}
\end{corollary}

\begin{proof}
Part (i) is exactly \eqref{eq:defect-coordinate-flow}. Iterating (i) gives
\[
L(T_t^{\,n}(p))=(1-t)^nL(p).
\]
Since
\[
L(T_{1-(1-t)^n}(p))=(1-(1-(1-t)^n))L(p)=(1-t)^nL(p),
\]
injectivity of $L$ proves \eqref{eq:iterate-formula}. If $t\neq 0$ and $T_t(p)=p$, then $(1-t)L(p)=L(p)$, hence $L(p)=0$ and therefore $p=\phi^{-1}$. This proves uniqueness of the fixed point. Finally, if $0<t<2$, then $|1-t|<1$, and \eqref{eq:global-attractor} follows immediately from the iterated formula for $L(T_t^{\,n}(p))$.
\end{proof}

The semigroup structure extends from the parameter space to the free-energy level. Since the Perron eigenvalue $\rho_p(t)$ satisfies a quadratic whose coefficients are entrywise powers of the transition probabilities, composing two tilts produces a factorization of characteristic polynomials that yields a global identity for the scaled cumulant generating function.

\begin{theorem}[Free-energy composition law]\label{thm:composition}
For the one-step Markov family on $X_\phi$, the scaled cumulant generating functions satisfy
\begin{equation}\label{eq:composition}
F_{q_p(t)}(s) = F_p(s+t-st) - (1-s)\,F_p(t)
\end{equation}
for all $p\in(0,1)$ and all $s,t\in\RR$. Consequently:
\begin{enumerate}[label=(\roman*),leftmargin=*,itemsep=2pt]
  \item The entropy rate of the tilted measure satisfies
  \begin{equation}\label{eq:entropy-composition}
  h(q_p(t)) = (1-t)\,F_p'(t) + F_p(t).
  \end{equation}
  \item The R\'enyi entropy rate
  \[
  R_\alpha(p):=\frac{1}{1-\alpha}\,F_p(1-\alpha), \qquad \alpha>0,\;\alpha\neq 1,
  \]
  satisfies $R_\alpha(\phi^{-1})=\log\phi$ for every $\alpha>0$, $\alpha\neq 1$. Since $R_1(\phi^{-1})=h(\phi^{-1})=\log\phi$ as well, the R\'enyi spectrum of the Parry measure is constant across all orders, and among one-step Markov laws on~$X_\phi$ this constancy characterizes the Parry measure.
\end{enumerate}
\end{theorem}

\begin{proof}
The Perron eigenvalue $\rho_{q_p(t)}(s)$ satisfies
\[
\lambda^2 - q_p(t)^{1-s}\,\lambda - (1-q_p(t))^{1-s}=0.
\]
Substituting $q_p(t)=p^{1-t}/\rho_p(t)$ and $1-q_p(t)=(1-p)^{1-t}/\rho_p(t)^2$ from Theorem~\ref{thm:tilted-free-energy}(ii) and setting $\Lambda:=\rho_p(t)^{1-s}\,\rho_{q_p(t)}(s)$, one obtains
\[
\Lambda^2 - p^{(1-s)(1-t)}\Lambda - (1-p)^{(1-s)(1-t)}=0.
\]
Since $(1-s)(1-t)=1-(s+t-st)$, this is the characteristic equation for $\rho_p(s+t-st)$. Both $\Lambda$ and $\rho_p(s+t-st)$ are positive roots of the same quadratic, hence
\[
\rho_p(t)^{1-s}\,\rho_{q_p(t)}(s)=\rho_p(s+t-st).
\]
Taking logarithms yields \eqref{eq:composition}. Part~(i) follows by differentiating at $s=0$ and using $F_q'(0)=h(q)$. For~(ii), set $t=1$ in \eqref{eq:composition}: since $q_p(1)=\phi^{-1}$ and $s+1-s=1$,
\[
F_{\phi^{-1}}(s) = F_p(1)-(1-s)\,F_p(1)=s\log\phi,
\]
so $R_\alpha(\phi^{-1})=F_{\phi^{-1}}(1-\alpha)/(1-\alpha)=(1-\alpha)\log\phi/(1-\alpha)=\log\phi$. Conversely, if $R_\alpha(p)$ is constant in $\alpha$, then in particular $\lim_{\alpha\to 1}R_\alpha(p)=h(p)$ equals $R_0(p)=F_p(1)=\log\phi$, so $h(p)=\log\phi$ and $p=\phi^{-1}$.
\end{proof}

\begin{corollary}[Variance transfer under tilts]\label{cor:variance-transfer}
The asymptotic jitter of the tilted measure satisfies
\begin{equation}\label{eq:variance-transfer}
\sigma^2(q_p(t)) = (1-t)^2\,F_p''(t).
\end{equation}
Equivalently, using Proposition~\ref{prop:variance-formula} and $L(q_p(t))=(1-t)L(p)$,
\begin{equation}\label{eq:Fpp-explicit}
F_p''(t) = \frac{q_p(t)\,(1-q_p(t))}{(2-q_p(t))^3}\,L(p)^2.
\end{equation}
\end{corollary}

\begin{proof}
Differentiating \eqref{eq:composition} twice with respect to $s$ gives $F_{q_p(t)}''(s)=(1-t)^2\,F_p''(s+t-st)$. Setting $s=0$ and using $F_q''(0)=\sigma^2(q)$ yields \eqref{eq:variance-transfer}. For \eqref{eq:Fpp-explicit}, substitute the explicit formula $\sigma^2(q)=q(1-q)(2-q)^{-3}L(q)^2$ with $q=q_p(t)$ and $L(q_p(t))=(1-t)L(p)$ into \eqref{eq:variance-transfer}, and cancel $(1-t)^2$.
\end{proof}

Since $F_p$ is finite and real analytic on all of~$\RR$, the G\"artner--Ellis theorem applies. The identity \eqref{eq:relative-entropy-along-tilt} is a manifestation of the standard Legendre duality between the scaled cumulant generating function and the rate function; the composition law of Theorem~\ref{thm:composition} makes this duality fully explicit by expressing the free energy of the tilted measure in terms of the free energy of the original.

\begin{corollary}[Large deviations in $q$-coordinates]\label{cor:tilted-ldp}
Let
\[
\Lambda_p(t):=F_p(t)-t h(p).
\]
Then the centered cylinder information satisfies a large-deviation principle:
\[
\frac{\tau_n-nh(p)}{n}
\]
obeys a good large-deviation principle under $\mu_p$ with rate function
\[
I_p(a):=\sup_{t\in\RR}\{ta-\Lambda_p(t)\}.
\]
Moreover, on the parametrized branch $a=\Lambda_p'(t)$ one has
\begin{equation}\label{eq:parametric-rate}
I_p(\Lambda_p'(t))=D(\mu_{q_p(t)}\|\mu_p).
\end{equation}
If $p\neq \phi^{-1}$, then for every $q\in(0,1)$ one has
\begin{equation}\label{eq:q-coordinate-rate}
I_p\bigl(h(q)-h(p)+D(\mu_q\|\mu_p)\bigr)=D(\mu_q\|\mu_p).
\end{equation}
In particular,
\begin{align}
\Lambda_p(1)&=\log\phi-h(p)=D(\mu_p\|\mu_\phi), \label{eq:lambda1-kl}\\
\Lambda_p'(1)&=D(\mu_p\|\mu_\phi)+D(\mu_\phi\|\mu_p). \label{eq:lambda1-prime-jeffreys}
\end{align}
\end{corollary}

\begin{proof}
By Theorem~\ref{thm:tilted-free-energy}, $F_p$ is finite and real analytic on all of $\RR$, hence the same is true for $\Lambda_p$. The large-deviation principle therefore follows from the G\"artner--Ellis theorem \cite{DemboZeitouni2010}. Since $I_p(\Lambda_p'(t))=t\Lambda_p'(t)-\Lambda_p(t)$ on the differentiability branch, \eqref{eq:parametric-rate} is exactly the identity
\[
t\Lambda_p'(t)-\Lambda_p(t)=tF_p'(t)-F_p(t),
\]
together with \eqref{eq:relative-entropy-along-tilt}.

If $p\neq \phi^{-1}$, then $L(p)\neq 0$. Since $L:(0,1)\to\RR$ is bijective, for every $q\in(0,1)$ there exists a unique
\[
t(q):=1-\frac{L(q)}{L(p)}
\]
such that $q=q_p(t(q))$. Substituting $t=t(q)$ into \eqref{eq:parametric-rate} gives \eqref{eq:q-coordinate-rate}.

At $t=1$, Theorem~\ref{thm:tilted-free-energy} gives $q_p(1)=\phi^{-1}$ and
\[
F_p(1)=\log\rho_p(1)=\log\phi.
\]
Hence
\[
\Lambda_p(1)=\log\phi-h(p).
\]
Using \eqref{eq:parametric-rate} with $t=1$ yields
\[
I_p(\Lambda_p'(1))=D(\mu_\phi\|\mu_p).
\]
Moreover,
\[
\begin{aligned}
D(\mu_p\|\mu_\phi)
&=
\frac{1}{2-p}
\left[
p\log\frac{p}{\phi^{-1}}
+(1-p)\log\frac{1-p}{\phi^{-2}}
\right] \\
&=
\frac{p\log p+(1-p)\log(1-p)}{2-p}+\log\phi
=
\log\phi-h(p)
=
\Lambda_p(1),
\end{aligned}
\]
which is \eqref{eq:lambda1-kl}. Finally, \eqref{eq:Fprime-cross-entropy} with $t=1$ gives
\[
F_p'(1)=h(\phi^{-1})+D(\mu_\phi\|\mu_p)=\log\phi+D(\mu_\phi\|\mu_p),
\]
and subtracting $h(p)$ yields \eqref{eq:lambda1-prime-jeffreys}.
\end{proof}

The following identity connects the jitter directly to the derivative of the entropy deficit $\Delta h(p)=\log\phi-h(p)$, and equivalently to the derivative of the KL divergence to the Parry measure.

\begin{proposition}[Jitter and entropy deficit]\label{prop:differential-identity}
For the one-step Markov family on $X_\phi$,
\begin{equation}\label{eq:differential-identity}
\sigma^2(p)=p(1-p)(2-p)\,\bigl(\Delta h'(p)\bigr)^2,
\qquad
\Delta h(p):=\log\phi-h(p).
\end{equation}
Equivalently, using \eqref{eq:lambda1-kl},
\begin{equation}\label{eq:differential-identity-kl}
\sigma^2(p)=p(1-p)(2-p)\left(\frac{d}{dp}D(\mu_p\|\mu_\phi)\right)^2.
\end{equation}
\end{proposition}

\begin{proof}
From the computation
\[
h'(p)=\frac{L(p)}{(2-p)^2},
\qquad
L(p):=\log\frac{1-p}{p^2},
\]
we have $\Delta h'(p)=-h'(p)$, so
\[
\bigl(\Delta h'(p)\bigr)^2=\frac{L(p)^2}{(2-p)^4}.
\]
Multiplying by $p(1-p)(2-p)$ and using Proposition~\ref{prop:variance-formula} gives
\[
p(1-p)(2-p)\,\bigl(\Delta h'(p)\bigr)^2
=
\frac{p(1-p)}{(2-p)^3}L(p)^2
=
\sigma^2(p).
\]
The equivalent form \eqref{eq:differential-identity-kl} follows from \eqref{eq:lambda1-kl}.
\end{proof}

\begin{remark}[Fisher information]\label{rem:fisher}
The Fisher information rate of the one-step family $\{\mu_p\}$ is
\[
I(p)
:=
\EE_{\pi_p}\!\left[\left(\frac{\partial}{\partial p}\log P_p(X_0,X_1)\right)^{\!2}\right]
=
\frac{1}{p(1-p)(2-p)}.
\]
Identity \eqref{eq:differential-identity} is therefore equivalent to
\begin{equation}\label{eq:fisher-identity}
\sigma^2(p)=\frac{\bigl(h'(p)\bigr)^2}{I(p)}.
\end{equation}
The vanishing $\sigma^2(\phi^{-1})=0$ reflects the fact that the Parry point is the unique critical point of $h$ on the one-step family: $h'(\phi^{-1})=0$. At the Parry point one has $h''(\phi^{-1})=-I(\phi^{-1})$, which is the standard relation between the entropy Hessian and the Fisher information at an entropy-maximizing parameter in exponential families \cite{Amari2016}. Combined with \eqref{eq:fisher-identity}, this directly yields the quadratic coefficient $2$ in Theorem~\ref{thm:second-order}: near $p_\phi:=\phi^{-1}$,
\[
\sigma^2(p)\approx I(p_\phi)(p-p_\phi)^2,
\qquad
\Delta h(p)\approx \tfrac{1}{2}I(p_\phi)(p-p_\phi)^2,
\]
so their ratio converges to $2$.
\end{remark}

\begin{figure}[ht]
\centering
\includegraphics[width=\textwidth]{fig_jitter_tilt.pdf}
\caption{(a) The asymptotic jitter $\sigma^2(p)$ (solid) and twice the entropy deficit $2\,\Delta h(p)$ (dashed) as functions of the Markov parameter $p$. Both vanish at the Parry point $p=\phi^{-1}$ and are tangent there, illustrating the quadratic comparison $\sigma^2\sim 2\,\Delta h$ of Theorem~\ref{thm:second-order}. (b) Tilt orbits $T_t^{\,n}(p_0)$ for $t=0.5$ and seven initial values $p_0$, converging to the Parry point $\phi^{-1}\approx 0.618$ as predicted by Corollary~\ref{cor:tilt-dynamics}(iii).}
\label{fig:jitter-tilt}
\end{figure}

\section{Tilt Closure for General Shifts of Finite Type}\label{sec:general-tilt}

The tilt closure property of Theorem~\ref{thm:tilted-free-energy}---that the Doob transform of the tilted transfer matrix is again a one-step Markov law on the same shift---is not specific to the golden-mean shift. In this section we show that it holds for every mixing shift of finite type and identify the structural property that permits the linearization of Corollary~\ref{cor:tilt-dynamics}.

\begin{proposition}[General tilt closure]\label{prop:general-tilt}
Let $(\Sigma_A,\sigma)$ be a topologically mixing shift of finite type with $k\times k$ adjacency matrix $A$. Let $P$ be a positive stochastic matrix compatible with $A$, and let $\mu_P$ be the corresponding stationary one-step Markov measure. For $t\in\RR$, define the tilted matrix
\[
B_P(t)_{ij}:=
\begin{cases}
P_{ij}^{1-t} & \text{if } A_{ij}=1,\\
0 & \text{otherwise.}
\end{cases}
\]
Then:
\begin{enumerate}[label=(\roman*),leftmargin=*,itemsep=2pt]
  \item $B_P(t)$ is a primitive nonnegative matrix with the same zero pattern as $A$.
  \item Its Doob transform $Q_t(i,j):=B_P(t)_{ij}\,r_t(j)/(\rho_t\,r_t(i))$, where $r_t$ is the positive right Perron eigenvector and $\rho_t$ the Perron eigenvalue, is a positive stochastic matrix compatible with $A$.
  \item The scaled cumulant generating function of $\tau_n$ under $\mu_P$ is $F_P(t)=\log\rho_t$.
\end{enumerate}
\end{proposition}

\begin{proof}
Since $P_{ij}>0$ whenever $A_{ij}=1$, the matrix $B_P(t)$ has the same positivity pattern as $A$ for every $t\in\RR$. Primitivity is inherited from the irreducibility and aperiodicity of $A$. For (ii), let $\rho_t$ and $r_t$ be the Perron eigenvalue and positive right eigenvector of $B_P(t)$. The Doob transform $Q_t(i,j):=B_P(t)_{ij}\,r_t(j)/(\rho_t\,r_t(i))$ satisfies $Q_t(i,j)=0$ whenever $A_{ij}=0$ (since $B_P(t)_{ij}=0$), and the row sums equal~$1$ because
\[
\sum_j Q_t(i,j)=\frac{1}{\rho_t\,r_t(i)}\sum_j B_P(t)_{ij}\,r_t(j)=\frac{\rho_t\,r_t(i)}{\rho_t\,r_t(i)}=1.
\]
Part (iii) is the standard Perron--Frobenius argument for the moment generating function of a Markov additive functional; see, e.g., Seneta~\cite{Seneta2006} or Dembo and Zeitouni~\cite{DemboZeitouni2010}.
\end{proof}

The parameter space of one-step Markov measures on $\Sigma_A$ has dimension $|E|-k$, where $|E|=\sum_{i,j}A_{ij}$ is the number of allowed edges and $k$ the number of states. For the golden-mean shift, $|E|=3$ and $k=2$, giving dimension~$1$: the row $(P_{10},P_{11})=(1,0)$ is completely determined by $A_{11}=0$, leaving the single free parameter $p=P_{00}$.

This one-dimensionality is the structural reason behind the global linearization of Corollary~\ref{cor:tilt-dynamics}. Any smooth self-map of an interval admitting a multiplicative coordinate can be linearized; the content of Theorem~\ref{thm:tilted-free-energy}(ii) is that $L$ serves as such a coordinate globally.

Although linearization is special, the fluctuation structure of cylinder information admits a useful normal form on any mixing SFT. Let $\mathcal{M}_A$ denote the set of all positive stochastic matrices compatible with $A$, and let $P^*$ be the Parry transition matrix,
\begin{equation}\label{eq:parry-transition}
P^*_{ij}:=\frac{A_{ij}\,r_j}{\lambda\,r_i},
\end{equation}
with stationary distribution $\pi^*_i:=\ell_ir_i/(\ell^\top r)$.

\begin{theorem}[Parry-relative fluctuation normal form]\label{thm:general-normal-form}
Let $P\in\mathcal{M}_A$ with stationary distribution $\pi_P$ and entropy rate $h(P)$. Define the edge deviation
\[
\Xi_P(i,j):=\log\frac{P^*_{ij}}{P_{ij}}+D(\mu_P\|\mu_{P^*}),
\qquad A_{ij}=1,
\]
and the endpoint correction
\[
B_P(i,j):=\log\frac{r_i}{\lambda\,\pi_{P,i}\,r_j}+D(\mu_P\|\mu_{P^*}).
\]
Then for every $n\ge 1$ and every admissible word $x_0^{n-1}$,
\begin{equation}\label{eq:general-normal-form}
\tau_n^{(P)}(x)-nh(P)=B_P(x_0,x_{n-1})+\sum_{m=0}^{n-2}\Xi_P(x_m,x_{m+1}).
\end{equation}
Moreover, $\int\Xi_P\,d\mu_P=0$, so $\sigma_\tau^2(\mu_P)=\sigma^2(\Xi_P,\mu_P)$.
\end{theorem}

\begin{proof}
From $\tau_n^{(P)}=-\log\pi_{P,x_0}-\sum_{m=0}^{n-2}\log P_{x_mx_{m+1}}$ and the Parry identity $-\log P^*_{ij}=\log\lambda+\log r_i-\log r_j$ (which telescopes along any path), one obtains
\[
\tau_n^{(P)}=-\log\pi_{P,x_0}+(n-1)\log\lambda+\log r_{x_0}-\log r_{x_{n-1}}+\sum_{m=0}^{n-2}\log\frac{P^*_{x_mx_{m+1}}}{P_{x_mx_{m+1}}}.
\]
Since $\log\lambda-h(P)=D(\mu_P\|\mu_{P^*})$, grouping the boundary terms gives \eqref{eq:general-normal-form}. The identity $\int\Xi_P\,d\mu_P=0$ follows from $\sum_i\pi_{P,i}\sum_jP_{ij}\log(P^*_{ij}/P_{ij})=-D(\mu_P\|\mu_{P^*})$.
\end{proof}

\begin{theorem}[Parry rigidity in the Markov simplex]\label{thm:general-rigidity}
Among all $P\in\mathcal{M}_A$, the Parry transition matrix~$P^*$ is the unique zero-jitter matrix: $\sigma_\tau^2(\mu_P)=0$ if and only if $P=P^*$.
\end{theorem}

\begin{proof}
By Theorem~\ref{thm:general-normal-form}, $\sigma_\tau^2(\mu_P)=0$ if and only if $\Xi_P$ is cohomologous to a constant: $\Xi_P(i,j)=c+u(i)-u(j)$ for all allowed edges. This gives $P_{ij}=e^{-D-c}\,P^*_{ij}\,e^{u(j)-u(i)}$ where $D=D(\mu_P\|\mu_{P^*})$. Setting $v_j:=e^{u(j)}$ and summing over $j$, the row-sum condition yields $\sum_jP^*_{ij}v_j=e^{D+c}v_i$. Since $P^*$ is primitive with Perron eigenvalue~$1$, Perron--Frobenius forces $e^{D+c}=1$ and $v$ constant. Hence $P=P^*$.
\end{proof}

The tilt map $T_t$ does not linearize globally when $\dim\mathcal{M}_A>1$, but its derivative at the Parry point retains a simple structure.

\begin{proposition}[Tangential linearization]\label{prop:tangential}
Let $\{P_\varepsilon\}$ be a $C^1$ curve in $\mathcal{M}_A$ with $P_0=P^*$, and let $s(i,j):=\frac{d}{d\varepsilon}\log P_\varepsilon(i,j)|_{\varepsilon=0}$ be the score function. Then
\[
\left.\frac{d}{d\varepsilon}\log T_t(P_\varepsilon)(i,j)\right|_{\varepsilon=0}=(1-t)\,s(i,j)+\psi_t(j)-\psi_t(i)-c_t
\]
for a state function $\psi_t$ and a constant $c_t$ arising from the perturbation of the Perron eigenvector and eigenvalue. In the quotient of edge functions modulo coboundaries and constants, $DT_t|_{P^*}$ acts as multiplication by $(1-t)$.
\end{proposition}

\begin{proof}
The Doob transform of $B_{P_\varepsilon}(t)$ with Perron data $(\rho_t(\varepsilon),r_t(\varepsilon))$ is
\[
T_t(P_\varepsilon)_{ij}=\frac{P_\varepsilon(i,j)^{1-t}\,r_t(\varepsilon,j)}{\rho_t(\varepsilon)\,r_t(\varepsilon,i)},
\]
so $\log T_t(P_\varepsilon)_{ij}=(1-t)\log P_\varepsilon(i,j)+\log r_t(\varepsilon,j)-\log r_t(\varepsilon,i)-\log\rho_t(\varepsilon)$.
At $\varepsilon=0$ one has $P_0=P^*$, and the Perron eigenvector of $B_{P^*}(t)$ is $r_t(0)=(r_i^t)_i$ with eigenvalue $\rho_t(0)=\lambda^{1-t}$ (since $P^*_{ij}=(A_{ij}r_j)/(\lambda r_i)$ gives $B_{P^*}(t)_{ij}=(r_j/(\lambda r_i))^{1-t}$, and $\sum_jB_{P^*}(t)_{ij}r_j^t=\lambda^{1-t}r_i^t$). Differentiating with respect to $\varepsilon$ at $\varepsilon=0$:
\[
\left.\frac{d}{d\varepsilon}\log T_t(P_\varepsilon)_{ij}\right|_{\varepsilon=0}
=
(1-t)\,s(i,j)+\psi_t(j)-\psi_t(i)-c_t,
\]
where $\psi_t(j):=\frac{d}{d\varepsilon}\log r_t(\varepsilon,j)|_{\varepsilon=0}$ and $c_t:=\frac{d}{d\varepsilon}\log\rho_t(\varepsilon)|_{\varepsilon=0}$. Since $\psi_t(j)-\psi_t(i)$ is a coboundary and $c_t$ is a constant, the quotient class $[DT_t|_{P^*}(s)]$ equals $(1-t)[s]$.
\end{proof}

\begin{example}[Full 2-shift]\label{ex:full-shift}
On the full shift $\{0,1\}^\NN$, the one-step Markov family is parametrized by
\[
P=\begin{pmatrix}p & 1-p\\q & 1-q\end{pmatrix},
\qquad 0<p,q<1,
\]
a two-dimensional family. The tilted Doob transform yields new parameters
\[
p'(t)=\frac{p^{1-t}}{\rho_t},
\qquad
q'(t)=\frac{q^{1-t}\,r_t}{\rho_t},
\]
where $r_t=(\rho_t-(1-q)^{1-t})/q^{1-t}$ is determined by the Perron eigenvector. The tilt map $T_t:(p,q)\mapsto (p'(t),q'(t))$ is a genuine two-dimensional dynamical system, and no single scalar coordinate linearizes it simultaneously in both components. The zero-jitter condition selects the Parry measure $p=q=\tfrac{1}{2}$ by a coboundary argument analogous to Theorem~\ref{thm:parry-rigidity}.
\end{example}

\begin{example}[Reversible chains on the full 2-shift]\label{ex:reversible}
The reversible subfamily $q=1-p$ forms a one-dimensional family preserved by the tilt, since the symmetry $B_P(t)_{ij}=B_P(t)_{ji}$ is inherited from $P_{ij}=P_{ji}$. With constant Perron eigenvector $r=(1,1)^\top$ and eigenvalue $\rho_t=p^{1-t}+(1-p)^{1-t}$, the Doob transform gives $p'(t)=p^{1-t}/\rho_t$. In the log-odds coordinate $M(p):=\log\frac{1-p}{p}$,
\[
M(p'(t))=(1-t)\,M(p),
\]
an exact analogue of $L(q_p(t))=(1-t)L(p)$ on the golden-mean shift. The unique fixed point $p=\tfrac{1}{2}$ is the Parry measure, and for $0<t<2$ it is globally attracting.
\end{example}

These examples show that linearization is not an accident of golden-mean arithmetic but a consequence of having a one-dimensional tilt-invariant family. What distinguishes the golden-mean shift is that the one-dimensionality arises \emph{naturally} from the adjacency constraints---the row $(1,0)$ in the transition matrix is forced---without restricting to a subfamily.

\section{H\"older equilibrium states and the second-order comparison}

We now place the golden-mean calculations in the broader Bowen--Ruelle framework. The next theorem collects standard consequences of coboundary theory \cite{Bowen1975,Ruelle1978,ParryPollicott1990} and exponential decay of correlations for Gibbs states \cite{Baladi2000}; we include the proof to make the connection to the cylinder-information setting explicit.

\begin{theorem}[Zero jitter for H\"older equilibrium states]\label{thm:gibbs}
Let $(\Sigma_A,\sigma)$ be a one-sided topologically mixing shift of finite type, let $f:\Sigma_A\to\RR$ be H\"older, and let $\mu_f$ be the equilibrium state of $f$. Define
\[
\tau_f(n,x):=-\log\mu_f([x_0^{n-1}]).
\]
Then:
\begin{enumerate}[label=(\roman*),leftmargin=*,itemsep=2pt]
  \item there exists $C_f<\infty$ such that
  \[
  \bigl|\tau_f(n,x)-(nP(f)-S_nf(x))\bigr|\le C_f
  \]
  for every $x$ and every $n\ge 1$;
  \item the asymptotic jitter rate exists and equals the asymptotic variance of $P(f)-f$;
  \item $\sigma_\tau^2(\mu_f)=0$ if and only if $f$ is H\"older-cohomologous to a constant.
\end{enumerate}
Equivalently, within the H\"older equilibrium states of H\"older potentials on $\Sigma_A$, the Parry measure is the unique zero-jitter measure.
\end{theorem}

\begin{proof}
The Gibbs property gives $K\ge 1$ such that
\[
K^{-1}
\le
\frac{\mu_f([x_0^{n-1}])}{\exp(S_nf(x)-nP(f))}
\le
K
\]
for every $n$-cylinder. Taking negative logarithms proves (i).

Let $g:=P(f)-f$ and define
\[
\beta_n(x):=\tau_f(n,x)-S_ng(x).
\]
Then $|\beta_n(x)|\le C_f$ for every $x$ and every $n$. Let
\[
\widehat g:=g-\int g\,d\mu_f.
\]
Because $g$ is H\"older whenever $f$ is H\"older, and because $\mu_f$ is a Gibbs equilibrium state for a H\"older potential on a mixing shift of finite type, correlations of H\"older observables decay exponentially in the Bowen--Ruelle theory \cite{Bowen1975,Ruelle1978,ParryPollicott1990}. Hence the Green--Kubo series
\[
\sigma_g^2
:=
\Var_{\mu_f}(\widehat g)
+2\sum_{k=1}^\infty \mathrm{Cov}_{\mu_f}(\widehat g,\widehat g\circ\sigma^k)
\]
converges absolutely and
\[
\Var_{\mu_f}(S_ng)=n\sigma_g^2+O(1).
\]
Since $\tau_f(n,\cdot)=S_ng+\beta_n$, we have
\[
\begin{aligned}
\bigl|\Var_{\mu_f}(\tau_f(n,\cdot))-\Var_{\mu_f}(S_ng)\bigr|
&\le \Var_{\mu_f}(\beta_n)
+2\sqrt{\Var_{\mu_f}(S_ng)\Var_{\mu_f}(\beta_n)} \\
&\le C_f^2+2C_f\sqrt{\Var_{\mu_f}(S_ng)}.
\end{aligned}
\]
Dividing by $n$ and using $\Var_{\mu_f}(S_ng)=n\sigma_g^2+O(1)$ shows that
\[
\frac{1}{n}\Var_{\mu_f}(\tau_f(n,\cdot))-\frac{1}{n}\Var_{\mu_f}(S_ng)\longrightarrow 0.
\]
Therefore the asymptotic jitter rate exists and equals $\sigma_g^2$, namely the asymptotic variance of $P(f)-f$. This proves (ii).

For H\"older observables on a mixing shift of finite type, zero asymptotic variance is equivalent to cohomology to a constant via a H\"older transfer function; this is standard in Bowen--Ruelle thermodynamic formalism \cite{Bowen1975,Ruelle1978,ParryPollicott1990}. Thus $\sigma_\tau^2(\mu_f)=0$ if and only if
\[
g=c+u-u\circ\sigma
\]
for some H\"older $u$ and constant $c$, equivalently
\[
f=(P(f)-c)-u+u\circ\sigma,
\]
so $f$ is cohomologous to a constant.

Conversely, assume that $f=a+v-v\circ\sigma$ for some constant $a$ and H\"older $v$. Then
\[
h_\nu(\sigma)+\int f\,d\nu
=
a+h_\nu(\sigma)
\]
for every $\sigma$-invariant measure $\nu$, so the equilibrium state is precisely the measure of maximal entropy, namely the Parry measure $\mu_P$ of $\Sigma_A$. This proves the equivalent uniqueness statement for zero-jitter equilibrium states. For that measure, if $\lambda$ is the Perron eigenvalue of $A$ and $\ell,r$ are positive left and right Perron eigenvectors, then every legal cylinder satisfies
\[
\mu_P([x_0^{n-1}])=\frac{\ell_{x_0}r_{x_{n-1}}}{\ell^\top r}\lambda^{-(n-1)}.
\]
Hence
\[
\tau_f(n,x)=\tau_P(n,x)=(n-1)\log\lambda+\log(\ell^\top r)-\log(\ell_{x_0}r_{x_{n-1}}),
\]
so $\tau_f(n,\cdot)-(n-1)\log\lambda$ depends only on the endpoints and is uniformly bounded. Therefore $\Var_{\mu_f}(\tau_f(n,\cdot))=O(1)$ and $\sigma_\tau^2(\mu_f)=0$. This proves (iii).
\end{proof}

\begin{remark}
We work throughout on the one-sided shift, which is the natural setting for the transfer-operator approach to exponential decay of correlations. The results extend to the two-sided natural extension via standard arguments: the one-sided Gibbs property determines the two-sided measure uniquely, and the coboundary characterization of zero variance carries over unchanged; see Parry and Pollicott~\cite{ParryPollicott1990} for the two-sided formulation.
\end{remark}

Returning to the golden-mean shift, we examine the rate at which the jitter vanishes as the parameter approaches the Parry point.

\begin{theorem}[Quadratic comparison near the Parry point]\label{thm:second-order}
Let
\[
\Delta h(p):=\log\phi-h(p).
\]
As $p\to\phi^{-1}$ in the one-step Markov family on $X_\phi$,
\begin{equation}\label{eq:second-order}
\sigma^2(p)=2\,\Delta h(p)+O\bigl((p-\phi^{-1})^3\bigr).
\end{equation}
Equivalently,
\[
\lim_{p\to\phi^{-1}}\frac{\sigma^2(p)}{\log\phi-h(p)}=2.
\]
Using \eqref{eq:lambda1-kl}, this is equivalently
\[
\sigma^2(p)=2\,D(\mu_p\|\mu_\phi)+O\bigl((p-\phi^{-1})^3\bigr).
\]
\end{theorem}

\begin{proof}
Set
\[
L(p):=\log\frac{1-p}{p^2}.
\]
By Proposition~\ref{prop:variance-formula},
\[
\sigma^2(p)=\frac{p(1-p)}{(2-p)^3}L(p)^2.
\]
Direct differentiation of
\[
h(p)=\frac{-p\log p-(1-p)\log(1-p)}{2-p}
\]
gives
\[
h'(p)=\frac{L(p)}{(2-p)^2}.
\]
At the Parry point $p_\phi:=\phi^{-1}$, the identity $1-p_\phi=p_\phi^2$ implies $L(p_\phi)=0$, hence both $h'(p_\phi)$ and $\sigma^2(p_\phi)$ vanish.

Taylor expansion yields
\[
L(p)=L'(p_\phi)(p-p_\phi)+O\bigl((p-p_\phi)^2\bigr),
\]
and therefore
\[
\sigma^2(p)
=
\frac{p_\phi(1-p_\phi)}{(2-p_\phi)^3}L'(p_\phi)^2(p-p_\phi)^2
+O\bigl((p-p_\phi)^3\bigr).
\]
On the other hand,
\[
\Delta h(p)
=
\frac{1}{2}\Delta h''(p_\phi)(p-p_\phi)^2
+O\bigl((p-p_\phi)^3\bigr),
\]
and since $h'(p)=L(p)(2-p)^{-2}$,
\[
\Delta h''(p_\phi)=-h''(p_\phi)=-\frac{L'(p_\phi)}{(2-p_\phi)^2}.
\]
Hence
\[
\frac{\sigma^2(p)}{\Delta h(p)}
\to
-2\,\frac{p_\phi(1-p_\phi)}{2-p_\phi}L'(p_\phi).
\]
Now
\[
L'(p)=-\frac{1}{1-p}-\frac{2}{p},
\]
and the relation $1-p_\phi=p_\phi^2$ gives
\[
-L'(p_\phi)=\frac{1+2p_\phi}{p_\phi^2}.
\]
Therefore
\[
-2\,\frac{p_\phi(1-p_\phi)}{2-p_\phi}L'(p_\phi)
=
2\,\frac{p_\phi^3}{2-p_\phi}\cdot\frac{1+2p_\phi}{p_\phi^2}
=
2\,\frac{p_\phi(1+2p_\phi)}{2-p_\phi}
=
2,
\]
which proves \eqref{eq:second-order}.
\end{proof}

The quadratic coefficient~$2$ is not special to the golden-mean shift. The proof of universality requires a smooth-dependence result for finite-state Markov chains, which we formulate as a lemma.

\begin{lemma}[Smooth dependence of asymptotic variance]\label{lem:smooth-variance}
Let $\{Q_t\}_{|t|<\epsilon}$ be a $C^2$ curve of primitive stochastic matrices on a finite state space, with $Q_0$ irreducible, and let $\{g_t\}$ be a $C^2$ family of edge observables with $g_0=0$. Write $g_t=t\,\dot{g}+O(t^2)$. Then
\[
\sigma^2(g_t,Q_t)=t^2\,\sigma^2(\dot{g},Q_0)+O(t^3).
\]
If furthermore $\sum_jQ_{0,ij}\,\dot{g}(i,j)=0$ for every $i$, then $\sigma^2(\dot{g},Q_0)=\sum_i\pi_{Q_0,i}\sum_jQ_{0,ij}\,\dot{g}(i,j)^2$.
\end{lemma}

\begin{proof}
The Gordin decomposition (Lemma~\ref{lem:markov-zero-variance}) expresses the asymptotic variance of an edge observable~$g$ under a finite-state ergodic chain with transition~$Q$ as $\sigma^2(g,Q)=\sum_i\pi_{Q,i}\sum_jQ_{ij}\,m(i,j)^2$, where $m(i,j)=g(i,j)-h+v(j)-v(i)$, $h=\sum_i\pi_{Q,i}\bar{g}(i)$, $\bar{g}(i)=\sum_jQ_{ij}g(i,j)$, and $v$ solves the Poisson equation $(Q-I)v=h\mathbf{1}-\bar{g}$. On a finite state space with $k$ states, the Poisson solution is $v=(Q-I+\mathbf{1}\pi^\top)^{-1}(h\mathbf{1}-\bar{g})$, which depends analytically on the entries of~$Q$ and~$g$ as long as $Q$ remains irreducible. Since $g_t=O(t)$, we have $\bar{g}_t=O(t)$, $h_t=O(t)$, $v_t=O(t)$, $m_t=O(t)$, and therefore $\sigma^2(g_t,Q_t)=O(t^2)$. The coefficient of~$t^2$ is $\sigma^2(\dot{g},Q_0)$, computed with $h_0=\sum_i\pi_{Q_0,i}\overline{\dot{g}}(i)$ and $v_0$ solving $(Q_0-I)v_0=h_0\mathbf{1}-\overline{\dot{g}}$. If $\overline{\dot{g}}(i)=\sum_jQ_{0,ij}\,\dot{g}(i,j)=0$ for every~$i$, then $h_0=0$, so $v_0=0$ and $m_0(i,j)=\dot{g}(i,j)$, giving $\sigma^2(\dot{g},Q_0)=\sum_i\pi_{Q_0,i}\sum_jQ_{0,ij}\,\dot{g}(i,j)^2$.
\end{proof}

The following theorem, which combines the Parry-relative normal form of Theorem~\ref{thm:general-normal-form} with Lemma~\ref{lem:smooth-variance}, establishes the universality of the coefficient~$2$.

\begin{theorem}[Universal quadratic coefficient]\label{thm:universal}
Let $(\Sigma_A,\sigma)$ be a topologically mixing SFT with Parry transition matrix~$P^*$ and stationary distribution~$\pi^*$. Let $\{P_t\}_{|t|<\epsilon}$ be a $C^2$ curve in $\mathcal{M}_A$ with $P_0=P^*$, and let
\[
s(i,j):=\left.\frac{d}{dt}\log P_t(i,j)\right|_{t=0}
\]
be the score function. Define the Fisher-type quadratic form
\[
\mathcal{I}_A(s):=\sum_i\pi^*_i\sum_jP^*_{ij}\,s(i,j)^2.
\]
Then
\begin{align}
D(\mu_{P_t}\|\mu_{P^*})&=\tfrac{1}{2}\,\mathcal{I}_A(s)\,t^2+O(t^3),\label{eq:universal-kl}\\
\log\lambda-h(P_t)&=\tfrac{1}{2}\,\mathcal{I}_A(s)\,t^2+O(t^3),\label{eq:universal-deficit}\\
\sigma_\tau^2(\mu_{P_t})&=\mathcal{I}_A(s)\,t^2+O(t^3).\label{eq:universal-jitter}
\end{align}
In particular,
\begin{equation}\label{eq:universal-two}
\lim_{t\to 0}\frac{\sigma_\tau^2(\mu_{P_t})}{\log\lambda-h(P_t)}=2.
\end{equation}
\end{theorem}

\begin{proof}
For \eqref{eq:universal-kl}, write $\delta_t(i,j):=P_t(i,j)-P^*_{ij}=tP^*_{ij}s(i,j)+O(t^2)$. Taylor expansion of $x\log(x/a)$ at $x=a$ gives
\[
D(\mu_{P_t}\|\mu_{P^*})=\sum_i\pi_{t,i}\sum_j\frac{\delta_t(i,j)^2}{2P^*_{ij}}+O(t^3),
\]
where the first-order term vanishes because $\sum_j\delta_t(i,j)=0$ (row-sum preservation). Since $\pi_t=\pi^*+O(t)$ and $\delta_t=tP^*_{ij}s(i,j)+O(t^2)$, this gives \eqref{eq:universal-kl}. The identity $\log\lambda-h(P)=D(\mu_P\|\mu_{P^*})$ for one-step Markov measures yields \eqref{eq:universal-deficit}.

For \eqref{eq:universal-jitter}, Theorem~\ref{thm:general-normal-form} gives $\sigma_\tau^2(\mu_{P_t})=\sigma^2(\Xi_{P_t},\mu_{P_t})$, where $\Xi_{P_0}=0$ (since $P_0=P^*$) and $\Xi_{P_t}=-t\,s+O(t^2)$. Lemma~\ref{lem:smooth-variance}, applied with $Q_t=P_t$ and $g_t=\Xi_{P_t}$, gives
\[
\sigma_\tau^2(\mu_{P_t})=t^2\,\sigma^2(s,\mu_{P^*})+O(t^3).
\]
The row-sum constraint $\sum_jP^*_{ij}s(i,j)=0$ and the second part of Lemma~\ref{lem:smooth-variance} yield
\[
\sigma^2(s,\mu_{P^*})=\EE_{\mu_{P^*}}\bigl[s(X_0,X_1)^2\bigr]=\mathcal{I}_A(s).
\]
Combining \eqref{eq:universal-kl} and \eqref{eq:universal-jitter} yields \eqref{eq:universal-two}.
\end{proof}

\begin{remark}
In the language of Riemannian geometry, \eqref{eq:universal-kl}--\eqref{eq:universal-jitter} state that at the Parry point,
\[
\mathrm{Hess}_{P^*}(\log\lambda-h)=\tfrac{1}{2}\,\mathcal{I}_A,
\qquad
\mathrm{Hess}_{P^*}(\sigma_\tau^2)=\mathcal{I}_A,
\]
so the entropy deficit and the jitter induce the same quadratic form at $P^*$, differing only by the universal factor~$2$. Theorem~\ref{thm:second-order} is the restriction of this identity to the one-dimensional golden-mean family.
\end{remark}

\section{Open questions and concluding remarks}

We summarize the main features and point out two directions that remain unexplored.

The golden-mean shift serves as the testing ground for a broader theory. The full-parameter normal form (Theorem~\ref{thm:normal-form}) separates cylinder-information fluctuations into an occupation-number bulk term and an endpoint correction, explaining structurally why the Parry measure has bounded variance. The tilt dynamics linearize globally in the $L$-coordinate, and the free-energy composition law (Theorem~\ref{thm:composition}) lifts this semigroup structure to the full fluctuation spectrum, yielding a variance transfer formula and characterizing the Parry measure through the constancy of its R\'enyi entropy rate.

These golden-mean results are instances of general phenomena. The Parry-relative normal form (Theorem~\ref{thm:general-normal-form}) holds on any mixing SFT and reduces the jitter to the asymptotic variance of the centered edge deviation $\Xi_P$. The Parry measure is the unique zero-jitter law in the full Markov simplex (Theorem~\ref{thm:general-rigidity}), and the quadratic coefficient $\sigma^2_\tau\sim 2(\log\lambda-h)$ is universal (Theorem~\ref{thm:universal}): the entropy deficit and the jitter induce the same Fisher-type quadratic form at the Parry point, differing only by the factor~$2$.

\medskip
\noindent\textbf{Linearization in higher dimensions.}
For a mixing SFT with parameter dimension $d=|E|-k>1$, does the tilt map $T_t$ on~$\RR^d$ admit a coordinate system---possibly curvilinear---in which it linearizes? Proposition~\ref{prop:tangential} shows that the derivative at the Parry point acts as multiplication by $(1-t)$ modulo coboundaries, but a global linearization seems unlikely when $d>1$.

\medskip
\noindent\textbf{Free-energy composition law for general SFTs.}
The composition law $F_{q_p(t)}(s)=F_p(s+t-st)-(1-s)F_p(t)$ of Theorem~\ref{thm:composition} relies on the golden-mean shift's one-dimensional parametrization. Does an analogous identity hold---in some tensorial or matrix form---for the multi-parameter free energy of a general mixing SFT? The Perron eigenvalue factorization in the proof extends formally, but the absence of a global coordinate like~$L$ makes the statement less immediate.

\medskip
\noindent\textbf{Defect distribution for larger alphabets.}
For a $k$-state mixing SFT, the Parry-measure defect $\tau_n-(n-1)\log\lambda$ depends on both endpoints and takes up to $k^2$ values. The exact defect distribution and its transient oscillations should be expressible in terms of the full spectrum of the transition matrix, generalizing the role of $-\phi^{-2}$ in Corollary~\ref{cor:finite-time-law}. We expect the bounded-variance property to persist, but the explicit formulas become considerably more involved.

\backmatter

\bmhead{Acknowledgments}

We thank the participants of the symbolic dynamics seminar for helpful discussions.

\section*{Declarations}

\begin{itemize}
\item \textbf{Funding.} No funding was received for conducting this study.
\item \textbf{Conflict of Interest.} The authors declare that they have no conflict of interest.
\item \textbf{Data Availability.} All data supporting the findings of this study are contained within the article. The Python script for reproducing Figure~\ref{fig:jitter-tilt} is available as supplementary material.
\item \textbf{Author Contributions.} H.~Ma conceived the study, developed the theoretical framework including the tilt dynamics on the Markov parameter space, the information normal form decomposition, the free-energy composition law, and the extensions to general shifts of finite type, and designed the overall research direction. W.~Zhang wrote the manuscript, independently verified all mathematical proofs and explicit formulas, performed the numerical computations, and prepared the figures.
\end{itemize}

\begin{thebibliography}{20}

\bibitem{Amari2016}
Amari, S.:
\textit{Information Geometry and Its Applications}.
Applied Mathematical Sciences, vol.~194.
Springer, Tokyo (2016)

\bibitem{Baladi2000}
Baladi, V.:
\textit{Positive Transfer Operators and Decay of Correlations}.
Advanced Series in Nonlinear Dynamics, vol.~16.
World Scientific, Singapore (2000)

\bibitem{Bowen1975}
Bowen, R.:
\textit{Equilibrium States and the Ergodic Theory of Anosov Diffeomorphisms}.
Lecture Notes in Mathematics, vol.~470.
Springer, Berlin (1975)

\bibitem{ChazottesUgalde2005}
Chazottes, J.-R., Ugalde, E.:
\textit{Entropy estimation and fluctuations of hitting and recurrence times for Gibbsian sources}.
Discrete Contin. Dyn. Syst. Ser.~B \textbf{5}(3), 565--586 (2005)

\bibitem{ClimenhagaThompson2012}
Climenhaga, V., Thompson, D.J.:
\textit{Intrinsic ergodicity beyond specification: $\beta$-shifts, $S$-gap shifts, and their factors}.
Israel J. Math. \textbf{192}, 785--817 (2012)

\bibitem{DemboZeitouni2010}
Dembo, A., Zeitouni, O.:
\textit{Large Deviations Techniques and Applications}, 2nd edn.
Stochastic Modelling and Applied Probability, vol.~38.
Springer, Berlin (2010)

\bibitem{FrischTamuz2022}
Frisch, J., Tamuz, O.:
\textit{Characteristic Measures of Symbolic Dynamical Systems}.
Ergodic Theory Dynam. Systems \textbf{42}(5), 1655--1661 (2022)

\bibitem{Gordin1969}
Gordin, M.I.:
\textit{The Central Limit Theorem for Stationary Processes}.
Soviet Math. Dokl. \textbf{10}, 1174--1176 (1969)

\bibitem{Haydn2012}
Haydn, N.:
\textit{The Central Limit Theorem for Uniformly Strong Mixing Measures}.
Stoch. Dyn. \textbf{12}, 1250006 (2012)

\bibitem{HuLin2013}
Hu, W.-G., Lin, S.-S.:
\textit{The Natural Measure of a Symbolic Dynamical System}.
arXiv:1308.2996 (2013)

\bibitem{Kifer1990}
Kifer, Y.:
\textit{Large Deviations in Dynamical Systems and Stochastic Processes}.
Trans. Amer. Math. Soc. \textbf{321}(2), 505--524 (1990)

\bibitem{KontoyiannisVerdu2014}
Kontoyiannis, I., Verd\'u, S.:
\textit{Optimal Lossless Data Compression: Non-Asymptotics and Asymptotics}.
IEEE Trans. Inform. Theory \textbf{60}(2), 777--795 (2014)

\bibitem{LindMarcus1995}
Lind, D., Marcus, B.:
\textit{An Introduction to Symbolic Dynamics and Coding}.
Cambridge University Press, Cambridge (1995)

\bibitem{Parry1964}
Parry, W.:
\textit{Intrinsic Markov Chains}.
Trans. Amer. Math. Soc. \textbf{112}, 55--66 (1964)

\bibitem{ParryPollicott1990}
Parry, W., Pollicott, M.:
\textit{Zeta Functions and the Periodic Orbit Structure of Hyperbolic Dynamics}.
Ast\'erisque \textbf{187--188} (1990)

\bibitem{PollicottSharp2009}
Pollicott, M., Sharp, R.:
\textit{Large deviations, fluctuations and shrinking intervals}.
Comm. Math. Phys. \textbf{290}(1), 321--334 (2009)

\bibitem{Ruelle1978}
Ruelle, D.:
\textit{Thermodynamic Formalism}.
Addison-Wesley, Reading (1978)

\bibitem{Seneta2006}
Seneta, E.:
\textit{Non-negative Matrices and Markov Chains}.
Springer Series in Statistics, 2nd edn.
Springer, New York (2006)

\bibitem{Shields1996}
Shields, P.C.:
\textit{The Ergodic Theory of Discrete Sample Paths}.
Graduate Studies in Mathematics, vol.~13.
American Mathematical Society, Providence (1996)

\bibitem{Walters1982}
Walters, P.:
\textit{An Introduction to Ergodic Theory}.
Springer, New York (1982)

\end{thebibliography}

\end{document}
